\documentclass[a4paper,11pt,fleqn]{scrartcl}%fleqn pour aligner les equations à gauche
\usepackage{../../Styles/Cours}

\newcommand{\chnum}{Chapitre 15}
\newcommand{\chtitle}{Etudes énergétiques en mécanique}
\newcommand{\dtype}{Cours}


\begin{document}
	\maketitle

\section{Energie cinétique d'un système}
    \subsection{Définition de l'énergie cinétique}
    \begin{cours}
        Dans un référentiel donné, l'énergie cinétique $E_c$ d'un système s'exprime par la relation : $$E_c = \dfrac{1}{2} m\cdot v^2 $$
        avec $E_c$ : l'énergie cinétique en joule (\unit{\joule})\\
        $m$ : la masse du système en kilogramme (\unit{\kilo\gram})\\
        $v$ : la vitesse du système en mètre par seconde (\unit{\meter\per\second})
    \end{cours}

    \subsection{Travail d'une force}
    Le travail d'une force est une grandeur physique permettant d'évaluer l'effet de cette force sur l'énergie cinétique d'un système au cours d'un mouvement.\\
    Le travail noté $W_{AB}(\vv{F})$ d'une force $\vv{F}$ dont le point d'application se déplace d'un point $A$ vers un point $B$ s'exprime par la relation scalaire : $$W_{AB}(\vv{F}) = \vv{F} \cdot \vv{AB} = F \cdot AB \cdot \cos(\alpha) $$
    avec : $W_{AB}(\vv{F})$ : le travail de la force $\vv{F}$ en joule (\unit{\joule})\\
    $F$ : a valeur de la force en newton (\unit{\newton})\\
    $AB$ : le déplacement en mètre (\unit{\meter})\\
    $\alpha$ : l'angle entre la direction de la force $\vv{F}$ et celle du déplacement $\vv{AB}$

    \begin{documentboxenv}{Cas 1}
        La force $\vv{F}$ ne travaille pas $\vv{F}$ et $\vv{AB}$ sont orthogonaux. $W_{AB}(\vv{F}) = \SI{0}{\joule}$
    \end{documentboxenv}

     \begin{documentboxenv}{cas 2}
        La force $\vv{F}$ travaille et ce travail est dit moteur car $$W_{AB}(\vv{F}) = \vv{F}\cdot \vv{AB}  = F \cdot AB \cdot \cos(0)> \SI{0}{\joule}$$
    \end{documentboxenv}

    \begin{documentboxenv}{cas 3}
        La force $\vv{F}$ travaille et ce travail est dit moteur car $$W_{AB}(\vv{F}) = \vv{F}\cdot \vv{AB}  = F \cdot AB \cdot \cos(0)> \SI{0}{\joule}$$
    \end{documentboxenv}
	
	
\end{document}